\section{Returning trajectories and collision boundaries}\label{sec:returns}

The escape windows leave an infinite family of returning trajectories.
Two further results restrict their possible brakes: a uniform angular-momentum
estimate at the first outer turn, and a local analysis of passages near
triple collision. Throughout this section, a collision of a positive-mass
trajectory ends its classical solution. Continuation through collision is
used only in the auxiliary regularized equations.

\subsection{The first outer turn}

On the returning side of the restricted separatrix, the first-turn time
$T(\phi)$ tends to infinity as $\phi\to0$. Since the binary mean anomaly
is $\phi+4\theta$, continuity supplies phases $\phi_n\to0$ with
$\phi_n+4T(\phi_n)=2\pi n$. At these phases the outer turn coincides
with a binary apocenter. Radial Kepler comparison gives
\[
 T_n={\pi n\over2}+o(1),\qquad Z_{{\rm t},n}=(2n)^{2/3}(1+o(1)).
\]
These are brakes of the regularized massless problem, so zeroth-order
convergence alone cannot exclude brakes of the full system.

The first transverse variation supplies the obstruction. On the centered
parabolic orbit of Section~\ref{sec:restricted}, let $p_-$ solve
\[
 p_-''={r^2-2z^2\over d^5}p_-,\qquad
 p_--z=O(|z|^{-1})\quad(\theta\to-\infty),\qquad
 d^2=z^2+r^2/4,
\]
where primes denote $\theta$ derivatives. Its Wronskian with $z$ satisfies
\[
 (zp_-'-z'p_-)'={3r^2zp_-\over2d^5},\qquad
 \mathcal W_\infty=\lim_{\theta\to\infty}(zp_-'-z'p_-).
\]
An outgoing solution $k_+$, normalized by $k_+/z\to1$ and
$zk_+'-z'k_+\to0$, gives by time reversal
$\mathcal W_\infty=2k_+(0)k_+'(0)$. Analytic Volterra bounds attach
this solution at finite radius: for $z\ge K\ge1$,
\[
 1\le{k_+\over z}\le Q_K={1\over1-1/(7K^2)},\qquad
 -{Q_K\over\sqrt{7/2}\,K^{3/2}}\le zk_+'-z'k_+\le0.
\]
Interval propagation of these bounds gives
\[
 k_+(0)>7/20,\qquad k_+'(0)>3/5,
 \qquad 11/25<\mathcal W_\infty<9/20.
\]

\begin{theorem}[Uniform first-turn exclusion]\label{thm:first-turn}
For exact tied trajectories with $B\to0$ whose restricted intercepts
approach the centered parabolic phase from the returning side, either a
prior classical collision occurs or, at the first outer turn,
\[
 {Y\times\dot Y\over B^{3/2}}\longrightarrow
 {\mathcal W_\infty\over2}>{21\over100}.
\]
The convergence is uniform in turn height. Hence no second brake occurs
at these first turns. On return to a large fixed scaled radius $K$, the
quotient differs from $\mathcal W_\infty/2$ by
$O(K^{-3/2})+o(1)$ as $B\to0$.
\end{theorem}
\begin{proof}
Set $\Lambda_B=(Z\times Z')/B$, $M=A+1$, and $N=M+B$.
The exact torque factorization is
\[
 \Lambda_B'=-{NA\over BM^2}(Z\times R)
 \bigl(|Z+R/M|^{-3}-|Z-AR/M|^{-3}\bigr).
\]
The squared denominators differ by
$2Z\cdot R+B^2|R|^2/M^2$. First-order incoming matching and the
divided LC normal equations retain $|Z\cdot R|\le CB|Z|$ on the
outgoing tail. A simultaneous binary-energy bootstrap keeps the LC
oscillator bounded. Writing $y=|Z|$, radial comparison gives
$-C/y^2\le y''\le-c/y^2$ and
$\int y^{-3}\,d\theta\le CK^{-3/2}$ from $y=K$ to any first turn.
This estimate includes the neighborhood where $y'=0$ and is uniform in
turn height. It bounds the change in $\Lambda_B$ by $CK^{-3/2}$.
At fixed $K$, the divided transverse matching gives
$\Lambda_B\to(zp_-'-z'p_-)/2$. Letting $K\to\infty$ proves the
claim. The same radial and torque estimates hold on the returning leg.
\end{proof}

For each fixed late restricted resonance, the transverse velocity is
also nonzero at small positive mass, but its coefficient tends to zero
as the resonance index increases. The angular-momentum estimate above
provides the uniform exclusion. It leaves the next central encounter
uncontrolled. At its centered parabolic endpoint, the asymptotic
coefficient $\gamma=\lim p_-/z$ satisfies
$-1/100<\gamma<-1/250$, and the outgoing Wronskian after the second encounter is
$-2\gamma\mathcal W_\infty>21/6250$; the interior returning phases
still require classification.

\subsection{The triple-collision endpoint and its joint scaling}

Another boundary of the returning region approaches equilateral triple
collision. The longitudinal and transverse unstable exponents in
logarithmic size are
\[
 \mu={1+\sqrt{19}\over4},\qquad
 \tau={1+\sqrt7\over4},\qquad 1<p:=\mu/\tau<2.
\]
Let $H_B,T_B$ be the corresponding normalized unstable amplitudes after
subtracting the local collision-stable graph. The odd transverse
coefficient satisfies $T_B/B\to\Theta_*\ne0$. Define
\[
 \rho_T=|T_B|^{1/\tau},\qquad
 \kappa_B={H_B\over|T_B|^p}.
\]
Scaling size by $\rho_T$ and time by $\rho_T^{3/2}$ produces a planar
massless limit indexed by $\kappa=\lim\kappa_B$. Its vector equations
are~\eqref{eq:restricted}, now with a zero-energy radial binary.
The longitudinal scaling gives the rectilinear exits when
$\kappa_B\to\pm\infty$.

The nonzero coefficient $\Theta_*$ is fixed by a certified transverse
intersection of the incoming parabolic curve with the collision-stable
curve. The latter has a unique connection in its normalized amplitude
bracket $[-0.24696,-0.24694]$. Along the terminal binary infall, the
canonical transverse collision modes have powers
$r^{\beta_\pm}$, $\beta_\pm=(3\pm\sqrt7)/4$, and Wronskian
$-\sqrt7$. An interval bound $W(J,P_+)<-1/2$ for the returned field
$J$ therefore proves its slow coefficient
$A_-=W(J,P_+)/(-\sqrt7)>0$. A second finite certificate, attached to
analytic incoming bounds, proves that the universal slow mode has positive
first-turn transfer $K_->0$. These signs exclude a brake near the second
turn on the selected local restricted branch.

For a finite light mass this conclusion needs a scale condition. If
$\varepsilon=(\pi-\phi)/4$ is the restricted center-to-collision time
gap and $\omega_0(B)\to0$ is the fixed-section state-matching error,
the transfer holds when
\[
 \omega_0(B)+B=o(\varepsilon^{(1+\sqrt{19})/6}).
\]
The transverse displacement relative to the inner size is then
$B\varepsilon^{-(1+\sqrt7)/6}=o(1)$. The complementary layer is
described by the planar $\kappa$ family. It cannot contain a brake
\emph{during} a sufficiently small near-triple passage: energy at a
brake requires the exact inequality
\[
 |R|>{A\over BU_0}
 ={1-B^2\over1+B^2-B^4}=1-2B^2+O(B^4).
\]
Possible brakes after re-expansion remain the issue.

\subsection{Planar collision boundaries and escape neighborhoods}

On a radial half of the limiting binary, set
$R=9^{1/3}t^{2/3}e_x$, $Z=9^{1/3}t^{2/3}w$, and $\zeta=\log t$.
Then
\[
 w_{\zeta\zeta}+\tfrac13w_\zeta=\nabla W(w),\qquad
 W(w)=\tfrac19\bigl(|w|^2+|w+e_x/2|^{-1}+|w-e_x/2|^{-1}\bigr).
\]
The shape energy obeys
$\mathcal E_\zeta=-|w_\zeta|^2/3$, where
$\mathcal E=|w_\zeta|^2/2-W$. For $d_\pm=|w\pm e_x/2|$,
\[
 9W=\sum_\pm(d_\pm^2/2+d_\pm^{-1})-1/4\ge11/4,
\]
with equality only at the equilateral shapes. A nonconstant orbit tending
to an equilateral shape has $\mathcal E>-11/36$ at finite times and
therefore cannot have zero shape velocity there.

Normalize the lower-equilateral stable family by
$w=(0,-\sqrt3/2)+P(a,b)$, where $P(a,b)=(a,b)+O_2$,
$a=-e^{-2\tau\zeta/3}$, and $b=\kappa e^{-2\mu\zeta/3}$.
These choices fix the scale and the meaning of $\kappa$ in the numerical
enclosures. Substitution in the shape equation determines its quartic
jet $P_4$ recursively. An exact homological inverse bound $30/17$ and a
contraction factor below $1/10$ give
$\|P-P_4\|_{1/200,1/2500}<1/125000000$ in the weighted coefficient
$\ell^1$ norm. Restriction to smaller polydiscs bounds both the state and
its $\kappa$ derivative. Higher jets improve the later collision enclosures.

In a selected light--primary chart $q=u^2$, $dt=|u|^2d\sigma$,
the equation $q_{tt}=-\Phi(q)+F(q,t)$ becomes
\[
 u_\sigma=v,\qquad
 v_\sigma=(h/2)u+(|u|^2\bar u/2)F,\qquad
 h_\sigma=2\operatorname{Re}(uv\bar F),
\]
with $2|v|^2-1-h|u|^2=0$. The other primary must remain separated.
The square collision map $(\kappa,\sigma)\mapsto(\Re u,\Im u)$
is tested by interval Newton. Four local roots have been enclosed:
\begin{center}
\small
\begin{tabular}{cll}
\toprule
Encounter chart & Primary & Enclosure for $\kappa$\\
\midrule
1 & positive & $[1.2679351752,1.2679351759]$\\
2 & negative & $[1.2640119191,1.2640119311]$\\
3 & positive & $[1.2640090993171269,1.2640090999643190]$\\
4 & negative & $[1.26400909893331527,1.26400909893331813]$\\
\bottomrule
\end{tabular}
\end{center}
Uniqueness is local to the parameter-time boxes in the certificates.
Earlier collisions are excluded for each listed root. At the first root,
$D_c=\det(\partial_\kappa u,v)\in(-3.52212,-3.38401)$; projection
onto $v^\perp$ gives
\[
 r_{\min,\rm local}=2D_c^2(\kappa-\kappa_c)^2
 +O(|\kappa-\kappa_c|^3).
\]
For positive mass, a colliding-pair center coordinate cancels the singular
force from the complementary equation. This gives a common analytic LC
chart, with $|v|^2=(1+B)/2$ at the first selected collision. Differentiated
inclination through the logarithmic saddle dwell then supplies a local
collision graph $\kappa_c(B,T,S,\Xi)$, where $S$ is the stable entrance
coordinate and $\Xi$ collects the remaining energy and clock deviations;
$\widehat\Xi$ denotes their Newtonian-scaled values.
The graph converges to the limiting root when
\[
 B|\log|T||+B+|T|^{2-p}+e^{-cL_T}+\|\widehat\Xi\|\to0,
 \qquad L_T=-\tau^{-1}\log|T|.
\]
The differentiated contraction uses a fixed weight
$\mu/2<\omega<\tau$; thus the quadratic Green denominator
$2\omega-\mu$ stays positive. This is the step that justifies derivative
convergence through the increasingly long saddle passage.

Finite regularized chart chains classify neighborhoods of the first,
second, and fourth roots as collision or escape. For example, ten adjacent
interval propagations cover
$[1.2679251755,1.2679451755]$ around the first root. They reach an
outgoing section with the strict margin
\[
 \dot\rho_0-{2\over(3/2)(\rho_0-r_0)}
 -\sqrt{4/r_0+2/100}-{3\over2}>0.65.
\]
This is Lemma~\ref{lem:hierarchical} with
$\varepsilon=1/100$, $c=3/2$, and limiting inner energy zero.
The scaled tidal allowance is $O(B)$, so the strict margin transfers to
sufficiently small compatible positive-mass states. The second-root
limiting box $[1.2640099161,1.2640126461]$ has the same local
finite-mass conclusion. The fourth-root chain requires simultaneous
two-centre regularization to resolve subsequent close passages; its
finite-mass-form escape margin exceeds $10.39$. A separate 131-tile
limiting cover proves collision or escape on
$[1.264009099014,1.264009099457]$; that explicit interval is a result
for the massless family. None of these local results classifies every
component of the $\kappa$ family.

\subsection{Exact collisions on the tied curve}

\begin{theorem}[Real collision parameters]
Infinitely many distinct real tied parameters $B_j\downarrow0$ have a
finite-time classical collision.
\end{theorem}
\begin{proof}
Each sufficiently late winding of the reference phase gives a disjoint
parameter bracket crossing the two sides of the certified triple endpoint.
An earlier collision anywhere in a bracket already supplies the assertion.
Otherwise, incoming matching transfers opposite signs of $H_B$ to its
endpoints. Within that bracket choose a segment from $H_B=0$ to
$H_B=B^{p/2}$ on which $0\le H_B\le B^{p/2}$. Local uniqueness of the
restricted intersection localizes this shrinking segment at the endpoint,
uniformly giving $T_B/B\to\Theta_*\ne0$ and the collision-graph
hypotheses. The continuous function
\[
 G(B)=H_B-|T_B|^p\kappa_c(B,T_B,S_B,\Xi_B)
\]
is negative at the first endpoint and positive at the second, since
$|T_B|^p=O(B^p)=o(B^{p/2})$. A zero is an exact light--heavy collision.
The disjoint brackets accumulate at zero.
\end{proof}
The roots are not known to be rational. On a bracket without an earlier
collision, the same normalized coordinate $H_B/|T_B|^p$ traverses the
first collision-or-escape interval. Its interior preimage supplies an
additional nonempty open tied window of nonperiodicity. Classification
of all remaining returning components, including their later turns,
is still open.
